\documentclass[10pt,twocolumn,letterpaper]{article}

% TEMPORARY CONTENT SCAFFOLD.
% Replace this preamble with the official CVPR 2027 author kit only after
% the venue template is verified and frozen.
\usepackage{times}
\usepackage{graphicx}
\usepackage{amsmath}
\usepackage{booktabs}
\usepackage{multirow}
\usepackage{microtype}
\usepackage{xcolor}

\title{When Should SAM 3 Write to Memory?\\
A Controlled Study of Signal-Informed Memory Admission for
Video Object Segmentation}

\author{Anonymous CVPR 2027 Submission}

\begin{document}
\maketitle

\begin{abstract}
Video object segmentation systems rely on memory to preserve object state
over time, yet not every prediction is equally suitable for future memory.
We study a controlled question: what information should a memory-write gate
use---quality signals, temporal signals, or identity signals?
We introduce a lightweight memory-admission mechanism for a frozen SAM 3
VOS/PVS tracker that physically accepts or blocks candidate memory writes
without fine-tuning the segmentation model.
We evaluate a nested family of manual, learned quality-temporal,
self-identity, and competitor-relative identity policies using
leakage-controlled cohorts, matched write rates, and video-clustered
inference.
On the frozen Hard TEST cohort, the learned quality-temporal gate improves
POR@30 over a write-rate-matched manual policy by 2.96 percentage points,
with a video-clustered BCa 95\% confidence interval of
[0.61, 6.43] percentage points.
In contrast, the tested native SAM 3 pointer identity signals do not
establish an incremental recovery benefit under the controlled comparisons.
Exact-budget neutral-selector controls further show that a lower write
frequency alone is insufficient to justify causal credit for a gating
policy.
These results support learned quality and temporal evidence as the
strongest demonstrated basis for memory admission in the tested setting,
while showing that write-rate control is essential for evaluating
memory-governance methods.
\end{abstract}

\section{Introduction}
\label{sec:intro}

Video object segmentation trackers increasingly rely on internal memory to
propagate object state across time.
A memory write is therefore not merely a bookkeeping operation: it changes
the information available to future predictions.
This creates a distinct control problem from frame-wise segmentation itself:
\emph{when should the tracker commit its current state to memory?}

Our study isolates that question in a frozen SAM 3 VOS/PVS substrate.
Rather than modifying the segmentation architecture or fine-tuning the
foundation model, we intervene only on candidate memory admission.
This design permits a controlled comparison of signal families while
preserving the underlying tracker.

The central research question is:
\emph{What information should a memory-write gate use---quality signals,
temporal signals, or identity signals?}

A key methodological difficulty is that different gates can write at
different rates.
A policy that writes less often can appear better or worse because of its
memory budget rather than because it selected more informative states.
We therefore treat realized write rate as an explicit comparison constraint
and complement matched-rate comparisons with exact-budget neutral controls.

Our strongest confirmatory result is that a learned quality-temporal policy
(B2) improves post-occlusion recovery over a matched manual policy (B1).
The tested native pointer-based identity extensions, however, do not
establish additional recovery benefit.
The result is deliberately narrower than claiming that gating universally
improves SAM 3 or that identity information is generally unhelpful.

Our contributions are:
\begin{itemize}
    \item a physically verified closed-loop memory-write intervention for
    frozen SAM 3 VOS/PVS;
    \item a nested evaluation of quality-temporal, self-identity, and
    competitor-relative identity signal families;
    \item a leakage-controlled evaluation protocol with matched write-rate
    requirements and video-clustered bootstrap inference;
    \item exact-budget neutral-selector controls that separate signal
    selection from write-frequency effects; and
    \item a bounded empirical finding that learned quality-temporal signals
    provide the strongest demonstrated basis for memory admission in the
    tested setting, while the tested native pointer identity signals provide
    no established incremental recovery benefit.
\end{itemize}

\section{Related Work}
\label{sec:related}

\paragraph{Selective memory and reliability.}
Prior VOS work already treats memory management as an explicit design
problem rather than an automatic consequence of propagation.
QDMN learns quality-aware dynamic memory admission in pre-foundation-model
VOS~\cite{QDMN}, while SAMURAI uses motion-aware memory selection in
SAM-family tracking~\cite{SAMURAI}.
DAM4SAM makes distractor-aware memory structure explicit~\cite{DAM4SAM},
and SENTRY validates candidate memory using training-free temporal and
trajectory reasoning~\cite{SENTRY}.
More recent SAM-generation studies further emphasize memory policy:
SAM3-DMS performs per-object reliability-based memory selection in a SAM 3
setting~\cite{SAM3-DMS}, and Rethinking Memory Design studies memory-policy
and memory-architecture choices across SAM-based trackers~\cite{RethinkingMemory}.
OAMVOS provides additional reliability- and occlusion-aware context for
selective memory updates~\cite{OAMVOS}.
These works make a broad \emph{memory filtering is new} claim untenable.
Our question is narrower: with the underlying SAM 3 tracker fixed, which
information is supported for the write-admission decision?

\paragraph{Identity-aware propagation and competitor reasoning.}
Identity information is well established as useful for associating objects
through time.
AOT, DeAOT, and AOST develop object-specific identity representations and
identity-aware propagation mechanisms for VOS~\cite{AOT,DeAOT,AOST}.
Recent systems also combine long-term memory, semantic identity, or
competitor-aware reasoning with modern segmentation models, including
ReMeDI-SAM3, SurgSLOT, CMR, and VOS-Agent
~\cite{ReMeDI-SAM3,SurgSLOT,CMR,VOS-Agent}.
CMR is especially relevant as a competitor-relative memory-readout neighbor,
but it does not implement the same write-admission intervention studied here.
Thus, prior identity-aware tracking motivates the scientific plausibility of
identity signals without establishing that native pointer identity adds
incremental value at write time once quality and temporal evidence are
already available.

\paragraph{Positioning of this study.}
The literature therefore supports all three candidate signal families:
quality, temporal consistency, and identity.
We do not compare unrelated end-to-end systems.
Instead, we hold the SAM 3 VOS/PVS substrate fixed and evaluate a nested
write-side ladder under a common leakage-controlled protocol, explicit
write-rate matching, exact-budget neutral controls, and video-clustered
inference.
This isolates the information available to memory admission rather than
claiming novelty for memory filtering itself.

\section{Controlled Memory-Write Admission}
\label{sec:method}

\subsection{Frozen SAM 3 substrate}

% TODO: describe build_sam3_video_model(), frozen weights, VOS/PVS memory
% path, and the intervention point using only repository-verified details.

\subsection{Closed-loop admission intervention}

% TODO: define candidate write, admit/block action, and verify that blocking
% changes future memory state while prediction continues.

\subsection{Signal-family ladder}

% TODO:
% B0: native write behavior.
% B1: manual quality/temporal gate.
% B2: learned quality/temporal gate.
% B3-S: B2 + self-identity information.
% B3-R: B3-S + competitor-relative identity information.
% B5: frozen comparison policy.
%
% Keep B3-R status and routing details consistent with canonical records.

\subsection{Gate constraints}

% TODO: <=50K trainable parameters, frozen SAM, FP32 identity cosine,
% no LoRA / no detector / no SAM fine-tuning.

\section{Evaluation Protocol}
\label{sec:protocol}

\subsection{Data and leakage boundaries}

% TODO: development-exposure boundary, Fresh DEV, Hard TEST,
% Representative TEST, outcome-independent selection, one-touch TEST.

\subsection{Endpoints}

% TODO: define POR@30 and ITR@30 exactly from the frozen endpoint record.

\subsection{Write-rate matching}

% TODO: explain frozen tolerance and why unmatched comparisons are
% descriptive rather than controlled causal contrasts.

\subsection{Statistical inference}

% TODO: video as statistical cluster; paired video-clustered BCa bootstrap;
% confidence level; effect-size reporting.

\subsection{Exact-budget neutral selectors}

% TODO: explain budget-preserving neutral timing control without overstating
% what the neutral comparison proves.

\section{Results}
\label{sec:results}

\subsection{Main Hard TEST result}

\begin{table}[t]
\centering
\caption{Hard TEST operating-point results. Values must be copied from the
frozen canonical result artifact before submission.}
\label{tab:hard-main}
\begin{tabular}{lccc}
\toprule
Policy & POR@30 & ITR@30 & Write rate \\
\midrule
B0   & 0.7411 & 0.0593 & 1.0000 \\
B1   & 0.7273 & 0.0514 & 0.3199 \\
B2   & 0.7569 & 0.0316 & 0.3233 \\
B3-S & 0.7648 & 0.0395 & 0.2906 \\
B3-R & 0.7648 & 0.0336 & 0.2923 \\
B5   & 0.7530 & 0.0336 & 0.3080 \\
\bottomrule
\end{tabular}
\end{table}

The primary matched comparison is B2 versus B1.
B2 improves POR@30 by 0.02964, with a video-clustered BCa 95\%
confidence interval of [0.00612, 0.06430].
Their realized write-rate difference is 0.00335, within the frozen matching
tolerance of 0.02.

\subsection{What do identity signals add?}

The properly write-rate-compatible B3-R versus B3-S comparison yields a
POR@30 difference of 0.00000, with a video-clustered BCa 95\% confidence
interval of [-0.01129, 0.01431].
Thus, the tested competitor-relative native pointer signal does not
establish an incremental recovery benefit.

The B3-S versus B2 comparison is not treated as a clean causal identity
contrast because its realized write-rate difference exceeds the frozen
matching tolerance.

\subsection{Why write-rate control matters}

% TODO: insert matched-rate table / effect-size figure.
% TODO: explain why B0 is descriptive only at the operating point.

\subsection{Exact-budget neutral controls}

% TODO: insert exact-budget signal-vs-neutral table.
% Required interpretation:
% Hard TEST signal-minus-neutral confidence intervals do not establish clear
% superiority of the signal-informed policies at the frozen operating point.

\subsection{Representative TEST}

% TODO: report the disjoint Representative TEST as external-validity evidence
% with its smaller event count stated clearly.

\section{Supplementary Post-Defense Full-Video Probe}
\label{sec:exp055}

EXP055 evaluates ordinary full-video segmentation quality on an
outcome-independently sampled 80-video supplementary cohort.
It is post-defense, nonconfirmatory evidence and does not modify the frozen
confirmatory TEST conclusions.

For pooled object-frame mean IoU, B0 obtains 0.73694 and B2 obtains 0.67810,
for a B2-minus-B0 difference of -0.05884.
The video-clustered BCa 95\% confidence interval is
[-0.11845, 0.00167].
The point estimate therefore favors B0, but the interval includes zero;
EXP055 establishes neither a statistically reliable full-video mIoU
improvement nor degradation for B2.

% TODO: likely move most EXP055 detail to supplement if main-page pressure
% requires it, retaining an explicit limitation/reference in the main paper.

\section{Discussion}
\label{sec:discussion}

The results support a narrow conclusion:
within the tested frozen SAM 3 substrate and evaluation protocol,
learned quality-temporal information provides the strongest demonstrated
basis for memory admission.

They do not support the broader claims that memory gating universally
improves SAM 3, that identity information is generally unhelpful, or that
the tested gate improves ordinary full-video segmentation quality.

% TODO: discuss non-monotonic write-rate sweep.
% TODO: discuss identity pointer availability and theft-event sparsity.
% TODO: discuss one-substrate / one-dataset limitation.
% TODO: discuss post-defense EXP055 cautiously.

\section{Conclusion}
\label{sec:conclusion}

Memory admission should be evaluated as a controlled decision problem rather
than as an incidental reduction in write frequency.
Under matched write rates, our learned quality-temporal gate improves
post-occlusion recovery over a manual policy, while the tested native SAM 3
identity signals provide no established incremental recovery benefit.
The exact-budget controls further show why memory-write policies require
budget-aware evaluation before gains can be attributed to their signals.

\bibliographystyle{plain}
\bibliography{references}

\end{document}
